\chapter{Flusso di elaborazione nel codice}
\label{chap:flusso-codice}

\section{Orchestrazione e sequenza generale}

\filepath{run.py} seleziona gli ID richiesti e delega ogni lavorazione a
\code{run_multiview_pipeline()} in \filepath{pipeline/orchestrator.py}. Da
questo punto l'orchestratore coordina l'intera trasformazione, senza contenere
direttamente gli algoritmi delle singole fasi.

Per ciascun nodo vengono costruite prima le osservazioni delle singole viste.
Gli artefatti VGGT, se abilitati, forniscono la geometria usata dalla ReID e
dalla successiva stima delle connessioni. Le istanze riconciliate diventano
quindi nodi di un grafo inizialmente privo di cicli; il refinement collega le
componenti rimaste separate, prima dell'annotazione dei diametri e della
serializzazione finale. La Figura~\ref{fig:pipeline-code-flow} riassume le
dipendenze tra queste fasi.

\begin{figure}[H]
  \centering
  \resizebox{\textwidth}{!}{\begin{tikzpicture}[
  node distance=5mm and 6mm,
  block/.style={draw=ManualBlue, rounded corners, fill=ManualLight,
                align=center, minimum height=10mm, text width=27mm,
                font=\small},
  artifact/.style={draw=ManualGreen, rounded corners, fill=white,
                   align=center, minimum height=9mm, text width=27mm,
                   font=\scriptsize},
  arrow/.style={-{Latex[length=2mm]}, thick, draw=ManualDark},
  artifactarrow/.style={-{Latex[length=2mm]}, dashed, draw=ManualGreen}
]
\node[block] (images) {Immagini\\multi-vista};
\node[block, right=of images] (observed) {Detection e\\grafi osservati};
\node[block, right=of observed] (artifacts) {Artefatti\\VGGT};
\node[block, right=of artifacts] (reid) {ReID\\multi-vista};
\node[artifact, below=of reid] (sceneout) {\code{scene_nodes.json}};

\node[block, below=15mm of artifacts] (hints) {Geometria e\\indizi 3D};
\node[block, left=of hints] (builder) {Costruzione\\diretta};
\node[block, left=of builder] (refine) {Refinement e\\connettività};
\node[block, left=of refine] (diam) {Annotazione\\diametri};
\node[artifact, below=of diam] (out) {\code{fused_graph.json}\\\code{bird_view.png}};

\draw[arrow] (images) -- (observed);
\draw[arrow] (observed) -- (artifacts);
\draw[arrow] (artifacts) -- (reid);
\draw[arrow] (reid) -- (hints);
\draw[artifactarrow] (reid) -- (sceneout);
\draw[arrow] (hints) -- (builder);
\draw[arrow] (builder) -- (refine);
\draw[arrow] (refine) -- (diam);
\draw[arrow] (diam) -- (out);
\end{tikzpicture}
}
  \caption{Le trasformazioni coordinate da \filepath{orchestrator.py}.}
  \label{fig:pipeline-code-flow}
\end{figure}

\section{Detection e riconciliazione delle osservazioni}

\subsection{Detection e grafi osservati}

\code{build_all_observed_graphs()} analizza ogni immagine tramite
\code{detect_and_decode_cached()}. Le detection sotto la soglia vengono
scartate, mentre i codici di diametro sono separati dai componenti topologici.
Ogni detection restante diventa un \code{ObservedNode} e le osservazioni della
stessa immagine vengono raccolte in un \code{ObservedLocalGraph}.

In questa fase non vengono stimati archi locali. La suddivisione per vista
conserva invece la provenienza delle osservazioni e introduce il primo vincolo
della ReID: due detection presenti nella stessa immagine non possono
rappresentare la stessa istanza fisica.

\projectfigure[0.80\textwidth]{detection_overlay.png}{Detection filtrate in una vista del nodo 11}{Output reale ottenuto dalla cache delle detection: in azzurro i componenti topologici e in arancione i codici di diametro.}

\subsection{Costruzione degli \code{SceneNode}}

\code{identify_scene_nodes()} in \filepath{node_reid.py} raggruppa prima le
osservazioni per codice. Per ogni codice costruisce un grafo di incompatibilità:
le osservazioni della stessa fotografia devono restare in gruppi diversi. La
colorazione greedy fornisce così un limite inferiore al numero di istanze
fisiche necessario a spiegare le co-occorrenze.

Quando la geometria è disponibile, il clustering agglomerativo parte dalle
singole osservazioni e fonde iterativamente la coppia compatibile migliore in
base alla distanza 3D tra le point cloud, senza unire detection della stessa
immagine. Il codice prevede anche un segnale basato sui vicinati osservati, ma
nella baseline corrente questo contributo resta inattivo per il limite descritto
nel Capitolo~\ref{chap:oggetti}. Senza geometria viene usata direttamente
l'assegnazione ottenuta dalla colorazione.

Ogni gruppo risultante diventa un \code{SceneNode}; l'associazione con le
osservazioni originali viene salvata in \filepath{scene_nodes.json}.

\projectfigure[0.76\textwidth]{reid_clusters.png}{Osservazioni raggruppate dalla ReID}{Tre istanze distinte con codice Shurikode 15; ogni riga raccoglie i crop attribuiti allo stesso \code{SceneNode} in cinque viste.}

\section{Geometria VGGT e costruzione degli indizi}

\subsection{Artefatti e geometria delle osservazioni}

Se VGGT è attivo per il nodo, l'orchestratore chiama
\code{ensure_vggt_artifacts()}. Per ogni vista viene riutilizzata la point map
già disponibile oppure, se consentito, viene generato il corrispondente file
\code{.npz}. All'interno delle bounding box vengono quindi campionati i punti
3D usati per stimare la geometria delle singole osservazioni.

La stessa geometria interviene in due momenti: supporta la riconciliazione
delle osservazioni durante la ReID e viene poi aggregata per \code{SceneNode}
per valutare le possibili connessioni.

\subsection{Indizi per le connessioni}

\code{build_vggt_edge_hints()} confronta le coppie di \code{SceneNode}
co-visibili. Per ciascuna raccoglie le viste di supporto e combina distanza tra
le point cloud o le porte, allineamento degli assi e ranking reciproco tra
vicini. Le coppie vietate vengono escluse, mentre le co-occorrenze note possono
rafforzare il punteggio.

Il risultato è ancora un insieme di \emph{hint}, non un insieme di archi. La
geometria aggregata e gli indizi vengono salvati in
\filepath{vggt_geometry.json} e possono essere riutilizzati quando la cache è
compatibile con gli \code{SceneNode} del run corrente.

\section{Costruzione, completamento e annotazione del grafo}

\subsection{Costruzione diretta}

Dopo il tentativo di riassorbire eventuali singleton spuri,
\code{build_direct_graph()} crea un \code{networkx.Graph} con un nodo per ogni
\code{SceneNode}. Le coppie suggerite da VGGT vengono ordinate per numero di
voti e punteggio e valutate in ordine greedy.

Una coppia viene accettata soltanto se non è vietata, se rispetta la valenza
disponibile dei componenti e se non introduce un ciclo. Il risultato iniziale
è quindi una foresta. Le coppie esaminate restano disponibili come
\code{BridgeEdge} per il completamento successivo.

\subsection{Completamento, diametri e output}

Se la foresta non è connessa, \code{force_spanning_tree()} prova prima i
\code{BridgeEdge} che uniscono componenti diverse. Se non bastano, applica un
fallback strutturale basato su co-occorrenze, triplette e supporto delle
istanze. Gli archi aggiunti in questa fase sono riconoscibili dall'attributo
\code{forced=true}.

\code{annotate_graph_diameters()} associa ogni codice diametro all'arco
visibile più vicino e propaga il valore lungo il gruppo di archi contiguo. La
propagazione si interrompe sui componenti configurati come possibili cambi di
diametro.

Infine \code{serialize_fused_graph()} salva struttura ed evidenze in
\filepath{fused_graph.json}; \code{render_bird_view()} produce la vista
d'insieme \filepath{bird_view.png}.
